% Fixed-H occupation transfer: valid first law; 1d flat, 3d CHM.
\documentclass[11pt]{article}
\usepackage[margin=1.05in]{geometry}
\usepackage{amsmath,amssymb,amsthm}
\usepackage[colorlinks=true,linkcolor=blue,citecolor=blue,urlcolor=blue]{hyperref}

\newtheorem*{admission}{Admission}

\title{\textbf{A Valid First-Law Probe at Fixed $H$:\\
1d Selects Enclosed Energy; 3d Balls Select CHM}}
\author{Robert W.\ McGwier\\
\small CTO, Cohere Technology Group\\
\small GitHub: \texttt{n4hy}}
\date{15 August 2026 --- Version 1}

\begin{document}
\maketitle

\begin{abstract}
\noindent
Paper~36 showed that a Hamiltonian point source is the wrong
variation: $\delta S\neq\mathrm{Tr}(K_{\mathrm{vac}}\Delta C)$.
This note never changes $H$. Occupation $\alpha$ is moved from
a localized occupied packet to its bipartite stagger.

The instrument passes. In 1d ($N=200$, $L=16$),
$\rho(\delta S,\mathrm{Tr}(K_{\mathrm{vac}}\Delta C))=0.957$,
but $P_{\mathrm{flat}}$ beats CHM ($R=0.060$ vs $0.266$;
$\rho_{\mathrm{flat}}=0.998$). Auditor: C$_{\mathrm{vac}}$
CONFIRMED, C2 REFUTED.

On $512$ balls of radius $2$ ($N=12$), the same construction
selects CHM: $R=0.018$ vs flat $0.102$ vs linear $0.029$.
C$_{\mathrm{vac}}=0.999$. Auditor $N=10$: C2 CONFIRMED.

A linear functional, yes. Einstein's equation, no.
The 1d theorem case still prefers Gauss, not the boost
weight. Not $8\pi G$. Not de~Sitter.
\end{abstract}

\section{What this is}

$K$ itself can be CHM-shaped (Paper~25) while the pairing of
this excitation against $K$ is the characteristic function of
a 1d interval and the CHM weight on a small 3d ball. That is
the first-law half of Faulkner--Guica--Hartman--Myers--Van
Raamsdonk as a lattice measurement of one variation. It is
not a derivation of $G_{\mu\nu}=8\pi GT_{\mu\nu}$.

\begin{admission}
I repaired the variation until the vacuum first law held. I
did not hide that 1d still selects enclosed energy, and I did
not write Einstein.
\end{admission}

\end{document}
